% Chapter 4, Figure 7(b): percent error in y_b, D4 engine (scan: figures/ch4/fig07b.png).
% Curves: digitized from the 1973 art by fig07b.py (the book gives neither the D4 thrust curve nor its
% propellant mass, so the interval solution (83)-(87) against eqs. (20), (21) and (27), (28) cannot be rerun);
% template fig06a.tex. Leader ends are points on the fig07b.csv curves.
\documentclass[11pt]{standalone}
\usepackage{tamrfig}
\begin{document}
\begin{tikzpicture}
  \begin{axis}[tamr,
      width=4.0in, height=2.5in,
      xmin=0.03, xmax=0.13, xtick={0.03,0.05,0.07,0.09,0.11,0.13}, minor xtick={0.04,0.06,0.08,0.10,0.12},
      xticklabel style={/pgf/number format/.cd, fixed, fixed zerofill, precision=2},
      ymin=-15, ymax=15, ytick distance=5,
      xlabel={$m_o$ (kg)}, ylabel={Percent error in $y_b$},
      legend pos=north east, legend style={nodes={inner sep=1pt}}]
    \draw[muted, line width=0.4pt] (axis cs:0.03,0) -- (axis cs:0.13,0);
    \addplot[series1] table[x=m0, y=fm_kmax, col sep=comma] {figures/v2/ch4/fig07b.csv};
    \addlegendentry{Fehskens-Malewicki}
    \addplot[series2] table[x=m0, y=cb_kmax, col sep=comma] {figures/v2/ch4/fig07b.csv};
    \addlegendentry{Caporaso-Bengen}
    \addplot[series1] table[x=m0, y=fm_kmin, col sep=comma] {figures/v2/ch4/fig07b.csv};
    \addplot[series2] table[x=m0, y=cb_kmin, col sep=comma] {figures/v2/ch4/fig07b.csv};
    % k labels with a leader to each method's curve for that k, placed as printed
    \node (kmax) at (axis cs:0.0696,-4.27) {\footnotesize $k_{\max}$};
    \draw[leader] (kmax.north) -- (axis cs:0.0683,-1.919);
    \draw[leader] (kmax.south) -- (axis cs:0.0648,-6.688);
    \node (kmin) at (axis cs:0.0426,-5.89) {\footnotesize $k_{\min}$};
    \draw[leader] (kmin.north) -- (axis cs:0.0422,-2.692);
    \draw[leader] (kmin.south) -- (axis cs:0.0403,-7.871);
  \end{axis}
\end{tikzpicture}
\end{document}
